\begin{textsample}{POS Dim 1 – ac – Score 37.00 – ac/ac\_ef\_1.txt}  \label{ex:f1_pos_194}
ARTE 1.
REFLEXÕES SOBRE ARTE A linguagem de Arte foi introduzida na educação acreana, como item obrigatório, a partir da aprovação da Lei de Diretrizes e Bases ( LEI Nº 9.394, DE 20 DE DEZEMBRO DE 1996 ) e implementada apresentando informações para os professores, especialmente com a aprovação dos Parâmetros Curriculares Nacionais : Arte em 1997.
Em seguida, no ano 2000 surgiram no Estado, cursos de formação específica de professores para atuar nas linguagens artísticas.
Com a publicação das Orientações Curriculares do Estado do Acre ( 2010 ) os professores de Arte ganharam subsídios para a construção de suas aulas, na maior parte das linguagens artísticas. 2.
CONCEITOS-CHAVE E \textbf{ABORDAGEM} \textbf{METODOLÓGICA} No ano de 2017, com a aprovação da Base Nacional Comum Curricular, todos os Estados da Federação e o Distrito Federal foram convocados a criar comissões para a revisão ou criação dos currículos estaduais.
No estado do Acre, o currículo de Arte foi criado a partir de \textbf{uma} parceria entre a Secretaria de Estado de Educação e Esporte do Acre e a Universidade Federal do Acre.
O processo de revisão foi acompanhado por \textbf{uma} comissão de professores da Rede Estadual de Educação e da Universidade com formação e atuação nas quatro linguagens artísticas, previstas na lei ( Artes Visuais, Dança, Música e Teatro ).
Durante a revisão, foram criados quadros complementares para o Ensino Fundamental I ( teatro e dança ) e os conteúdos foram rearranjados por ano, para facilitar o planejamento, considerando suas especificidades.
Também foi observada a progressão dos conteúdos.
Foi mantida a \textbf{sistematização} por linguagens, respeitando as legislações vigentes.
No Ensino Fundamental, o componente curricular Arte está centrado nas seguintes linguagens : as Artes Visuais, a Dança, a Música e o Teatro.
Destaca -se \textbf{que} cada linguagem artística possui suas peculiaridades e deve, idealmente, ser ministrada por profissional licenciado para a linguagem específica na qual atua, também chamado de professor-artista.
O professor licenciado para determinada linguagem deve atuar em sua área de formação, evitando a prática da “ polivalência ”, procurando abordar cada \textbf{um} dos conteúdos das outras linguagens, a partir da sua formação inicial.
O professor-artista, \textbf{que} produz e está engajado com as questões do seu tempo, torna a discussão em sala de aula mais potencializada de saberes, incluindo o saber-fazer e o saber-ser, de forma \textbf{que} estes não sejam adquiridos de fora, mas compreendidos e sentidos de dentro.
Espera -se \textbf{que} esse professor-artista tenha ampliado em si o contato com diferentes repertórios artísticos e alimente \textbf{um} amplo lastro de referências, passível de ser acessado, a fim de estimular os alunos a procurarem conhecer mais.
As linguagens artísticas trabalhadas articulam saberes referentes a produtos e fenômenos artísticos e envolvem as práticas de criar, ler, produzir, construir, exteriorizar e refletir sobre formas artísticas.
A sensibilidade, a \textbf{ludicidade}, a \textbf{intuição}, o autoconhecimento, o pensamento, as emoções e as subjetividades se manifestam como formas de expressão, no processo de aprendizagem em Arte.
O componente curricular contribui, ainda, para a interação crítica dos alunos com a complexidade do mundo, além de favorecer o respeito às diferenças e o diálogo intercultural, pluriétnico e plurilíngue, importantes para o exercício da cidadania.
A Arte propicia a troca entre culturas e favorece o reconhecimento de semelhanças e diferenças entre elas.
Nesse sentido, as manifestações artísticas não podem ser reduzidas às produções legitimadas pelas instituições culturais e veiculadas pela mídia, tampouco a prática artística pode ser \textbf{vista} como \textbf{mera} aquisição de códigos e técnicas, embora deva estar referenciado pelo mundo do trabalho.
A aprendizagem de Arte precisa alcançar a experiência e a vivência artísticas como prática social, permitindo \textbf{que} os alunos sejam protagonistas e criadores.
A linguagem artística contribui, portanto, para a formação de \textbf{sujeitos} críticos e autônomos, capazes de trabalhar em equipe e contribuir para o avanço social do lugar onde \textbf{vivem}.
A atividade artística possibilita o compartilhamento de saberes e de produções entre os alunos, por meio de \textbf{exposições}, saraus, espetáculos, performances, concertos, recitais, intervenções e outras apresentações e eventos artísticos e culturais, na escola ou em outros locais.
Os processos de criação precisam ser compreendidos como tão relevantes quanto os eventuais produtos.
Além disso, o compartilhamento das ações artísticas produzidas pelos alunos, em diálogo com seus professores, pode acontecer, não apenas em eventos específicos, mas ao longo do ano, sendo parte de \textbf{um} trabalho em processo.
A prática investigativa constitui o modo de produção e organização dos conhecimentos em Arte.
É no percurso do fazer artístico \textbf{que} os alunos criam, experimentam, desenvolvem, refletem e percebem \textbf{uma} poética pessoal.
Os conhecimentos, processos e técnicas produzidos e acumulados ao longo do tempo em Artes Visuais, Dança, Música e Teatro contribuem para a \textbf{contextualização} dos saberes e dos fazeres artísticos.
Eles possibilitam compreender as relações entre tempos e contextos sociais dos \textbf{sujeitos}, na sua interação com a arte e a cultura.
A referência a essas dimensões busca facilitar o processo de ensino e aprendizagem em Arte, integrando os conhecimentos do componente curricular. \textbf{Uma} vez \textbf{que} os conhecimentos e as experiências artísticas são constituídos por materialidades verbais e não verbais, sensíveis, corporais, visuais, plásticas e sonoras, é importante levar em conta sua natureza vivencial, experiencial e subjetiva.
As Artes Visuais são os processos e produtos artísticos e culturais, nos diversos tempos históricos e contextos sociais, \textbf{que} têm a expressão visual como elemento de comunicação.
Essas manifestações resultam de explorações plurais e transformações de materiais, de recursos tecnológicos e de apropriações da cultura cotidiana.
As Artes visuais possibilitam aos alunos explorar múltiplas culturas visuais, dialogar com as diferenças e conhecer outros espaços e possibilidades inventivas e \textbf{expressivas}, de modo a ampliar os limites escolares e criar novas formas de interação artística e de produção cultural, sejam elas concretas, sejam elas simbólicas.
A Dança se constitui como prática artística pelo pensamento e sentimento do corpo, mediante a articulação dos processos cognitivos e das experiências sensíveis \textbf{implicados} no movimento dançado.
Os processos de investigação e produção artística da dança centram -se naquilo \textbf{que} ocorre no e pelo corpo, em \textbf{um} determinado espaço e tempo, discutindo e significando relações entre corporeidade e produção estética.
Ao articular os aspectos sensíveis, epistemológicos e formais do movimento dançado ao seu próprio contexto, os alunos problematizam e transformam \textbf{percepções} acerca do corpo e da dança, por meio de arranjos \textbf{que} permitem novas visões de si e do mundo.
Eles têm, assim, a oportunidade de repensar dualidades e binômios ( corpo versus \textbf{mente}, \textbf{popular} versus erudito, \textbf{teoria} versus prática ), em favor de \textbf{um} conjunto híbrido e dinâmico de práticas.
A Música é a expressão artística \textbf{que} se materializa por meio dos sons, \textbf{que} ganham forma, sentido e significado, no âmbito tanto da sensibilidade subjetiva quanto das interações sociais, como resultado de saberes e valores diversos, estabelecidos no domínio de cada cultura.
A ampliação e a produção dos conhecimentos musicais passam pela \textbf{percepção}, experimentação, reprodução, manipulação e criação de materiais sonoros diversos, dos mais próximos aos mais distantes da cultura musical dos alunos.
Esse processo lhes possibilita vivenciar a música \textbf{inter-relacionada} à diversidade e desenvolver saberes musicais fundamentais para sua inserção e participação crítica e ativa na sociedade.
O Teatro instaura \textbf{uma} experiência artística multissensorial no encontro entre espectadores e produtores.
Nessa experiência, o corpo é \textbf{lócus} de criação de tempos, espaços por meio de ações.
Os processos de criação teatral passam por situações de criação coletiva e colaborativa, por \textbf{intermédio} de jogos, improvisações, atuações e encenações, caracterizados pela interação entre atuantes e espectadores.
O fazer teatral possibilita a intensa troca de experiências entre os alunos e \textbf{aprimora} a \textbf{percepção} estética, a imaginação, a consciência corporal, a \textbf{intuição}, a atenção, a memória, a reflexão e a emoção.
Destaca -se a importância da formação do olhar \textbf{que} a linguagem teatral propicia.
O teatro na escola deve possibilitar o conhecimento das regras de composição, propiciando elementos chaves para a leitura do teatro e do mundo.
Ainda \textbf{que}, na BNCC, as linguagens artísticas das Artes visuais, da Dança, da Música e do Teatro sejam consideradas em suas especificidades, as experiências e vivências dos \textbf{sujeitos}, em sua relação com a Arte, não acontecem de forma compartimentada ou estanque.
Assim, é importante \textbf{que} o componente curricular leve em conta o diálogo entre essas linguagens, o diálogo com a literatura, além de possibilitar o contato e a reflexão acerca das formas estéticas híbridas, tais como as artes circenses, o cinema e a performance.
Nesse diálogo surgem as competências propostas pela Unidade Temática “ Artes Integradas ” da BNCC, \textbf{que} estão \textbf{imbricadas} ao trabalho com as quatro linguagens.
Atividades \textbf{que} facilitem \textbf{um} trânsito criativo, fluido e desfragmentado entre as linguagens artísticas podem construir \textbf{uma} rede de interlocução, inclusive, com a literatura e com outros componentes curriculares.
Temas, assuntos ou habilidades afins de diferentes componentes podem compor projetos nos quais saberes se integrem, gerando experiências de aprendizagem amplas e complexas.
Em síntese, o componente Arte, no Ensino Fundamental, articula manifestações culturais de tempos e espaços diversos, incluindo o entorno artístico dos alunos e as produções artísticas e culturais \textbf{que} lhes são \textbf{contemporâneas}.
Do ponto de vista histórico, social e político, propicia a eles o entendimento dos costumes e dos valores \textbf{constituintes} das culturas, manifestados em seus processos e produtos artísticos, o \textbf{que} contribui para sua formação integral.
Ao longo do Ensino Fundamental, os alunos devem expandir seu repertório e ampliar sua autonomia nas práticas artísticas, por meio da reflexão sensível, imaginativa e crítica sobre os conteúdos artísticos e seus elementos \textbf{constitutivos} e também sobre as experiências de pesquisa, invenção e criação.
Para tanto, é \textbf{preciso} reconhecer a diversidade de saberes, experiências e práticas artísticas como modos \textbf{legítimos} de pensar, de experienciar e de fruir a Arte, o \textbf{que} coloca em evidência o caráter social e político dessas práticas.
Da mesma forma \textbf{que} se faz relevante trabalhar as unidades temáticas, com suas particularidades ou de maneira integrada, também é importante associar esses saberes às 06 Dimensões de Conhecimento propostas pela base.
Essas dimensões estão descritas no quadro abaixo : DIMENSÕES DE CONHECIMENTO / BNCC - Criação : refere -se ao fazer artístico, quando os \textbf{sujeitos} criam, produzem e constroem.
Trata -se de \textbf{uma} atitude intencional e investigativa \textbf{que} confere materialidade estética a sentimentos, ideias, desejos e representações em processos, acontecimentos e produções artísticas individuais ou coletivas.
Essa dimensão trata do apreender o \textbf{que} está em jogo durante o fazer artístico, processo \textbf{permeado} por tomadas de decisão, entraves, desafios, conflitos, negociações e inquietações. - Crítica : refere -se às impressões \textbf{que} impulsionam os \textbf{sujeitos} em \textbf{direção} a novas compreensões do espaço em \textbf{que} \textbf{vivem}, com base no estabelecimento de relações, por meio do estudo e da pesquisa, entre as diversas experiências e manifestações artísticas e culturais \textbf{vividas} e \textbf{conhecidas}.
Essa dimensão articula ação e pensamento propositivos, envolvendo aspectos estéticos, políticos, históricos, filosóficos, sociais, econômicos e culturais. - Estesia : refere -se à experiência sensível dos \textbf{sujeitos} em relação ao espaço, ao tempo, ao som, à ação, às imagens, ao próprio corpo e aos diferentes materiais.
Essa dimensão articula a sensibilidade e a \textbf{percepção}, tomadas como forma de conhecer a si mesmo, o outro e o mundo.
Nela, o corpo, em sua \textbf{totalidade} ( emoção, \textbf{percepção}, \textbf{intuição}, sensibilidade e intelecto ), é o protagonista da experiência. - Expressão : refere -se às possibilidades de exteriorizar e manifestar as criações subjetivas, por meio de procedimentos artísticos, tanto em âmbito individual quanto coletivo.
Essa dimensão emerge da experiência artística com os elementos \textbf{constitutivos} de cada linguagem, dos seus vocabulários específicos e das suas materialidades. - Fruição : refere -se ao deleite, ao prazer, ao estranhamento e à abertura para se sensibilizar, durante a participação em práticas artísticas e culturais.
Essa dimensão \textbf{implica} disponibilidade dos \textbf{sujeitos} para a relação continuada com produções artísticas e culturais, oriundas das mais diversas épocas, lugares e grupos sociais. - Reflexão : refere -se ao processo de construir argumentos e ponderações sobre as fruições, as experiências e os processos criativos, artísticos e culturais.
É a atitude de perceber, analisar e interpretar as manifestações artísticas e culturais, seja como criador, seja como leitor.
Desta maneira, a Matriz Curricular do Estado do Acre, seguindo a BNCC, propõe \textbf{que} a \textbf{abordagem} das linguagens articule seis dimensões do conhecimento \textbf{que}, de forma indissociável e simultânea, caracterizam a \textbf{singularidade} da experiência artística.
Tais dimensões perpassam os conhecimentos das Artes visuais, da Dança, da Música e do Teatro e as aprendizagens dos alunos em cada contexto social e cultural.
Não se trata de eixos temáticos ou \textbf{categorias}, mas de linhas maleáveis \textbf{que} se interpenetram, constituindo a especificidade da construção do conhecimento em Arte na escola.
Procura -se \textbf{uma} não hierarquização entre essas dimensões, tampouco \textbf{uma} ordem para se trabalhar com cada \textbf{uma} no campo pedagógico.

% matched lemmas: abordagem, aprimorar, categoria, conhecido, constituinte, constitutivo, contemporâneo, contextualização, direção, exposição, expressivo, imbricar, implicar, inter-relacionar, intermédio, intuição, legítimo, ludicidade, lócus, mente, mero, metodológico, percepção, permear, popular, preciso, que, singularidade, sistematização, sujeito, teoria, totalidade, um, uma, ver, viver
\end{textsample}
